\section{Evaluation}\label{sec:eval}
The detection and backtest numbers come from the SF=0.25 build with real weather (seed 42) of 20
October\prov{}, in the detection trade-off and backtest metric files under \texttt{docs/experiments}.

\subsection{Detection}
\texttt{make evaluate} scores every alert against the 35 injected outbreaks in \texttt{ground\_truth/}, the only
code allowed to read that folder. Consecutive alert weeks merge into one episode; an episode that overlaps no
outbreak window counts as a false alert. The sweep picks, for each method, the most sensitive setting that keeps
false alerts at or below two per district-year. Table~\ref{tab:detection} shows the chosen operating points.

\begin{table}[htbp]
\caption{Detection at the chosen operating points (SF=0.25)\prov}
\label{tab:detection}
\centering\small
\begin{tabular}{@{}lrr@{}}
\toprule
 & EARS C2 & CUSUM \\
\midrule
Setting & threshold 4.0 & $k=0.5$, $h=6$ \\
Outbreaks detected & 34 of 35 & 32 of 35 \\
Sensitivity & 0.97 & 0.91 \\
False alerts per district-year & 1.27 & 1.49 \\
Positive predictive value & 0.22 & 0.28 \\
Median days to detection & 8 & 9 \\
Detected within 14 days & 0.91 & 0.83 \\
\bottomrule
\end{tabular}
\end{table}

EARS C2 is the better default here: it finds more outbreaks, sooner, with fewer false alerts. We expected CUSUM
to do better on slow growth, but on the one slow outbreak both methods found it late, EARS C2 after 48 days and
CUSUM after 69.

\begin{teammate}{Samira (AN-2), Rafi (AN-3): confirm}
Add the C1 trade-off figure from \texttt{detection\_tradeoff.csv} (sensitivity against false alerts per
district-year, one curve per method, the chosen points marked), results by outbreak shape, and the artefact
facility with the naive detector and with ours. Replace with the v1.0 numbers.
\end{teammate}

\subsection{Forecast backtest}
The backtest refits the per-district NB-GLM at 52 rolling origins (every fourth origin, reused in between) and
forecasts one to four weeks ahead using only data up to each origin. The seasonal-naive baseline is scored on
exactly the same district, disease, origin and horizon cells. The weighted interval score~\cite{bracher2021wis}
uses eleven central intervals plus the median; relative WIS below 1 means the model beats the baseline.

\begin{table}[htbp]
\caption{Backtest over 52 origins, model against seasonal naive (SF=0.25)\prov}
\label{tab:backtest}
\centering\small
\setlength{\tabcolsep}{4pt}
\begin{tabular}{@{}llrrrr@{}}
\toprule
Disease & $h$ & MAE & MAE base & Cov.\,80\% & Rel.\ WIS \\
\midrule
Dengue & 1 & 39.4 & 42.5 & 0.83 & 0.88 \\
 & 4 & 41.4 & 44.9 & 0.82 & 0.89 \\
Diarrhoea & 1 & 18.5 & 26.3 & 0.80 & 0.70 \\
 & 4 & 19.3 & 26.4 & 0.80 & 0.74 \\
\midrule
All & 1--4 & 29.5 & 35.0 & 0.81 & 0.83 \\
\bottomrule
\end{tabular}
\end{table}

The model beats the baseline at every horizon for both diseases, by more for diarrhoea (relative WIS 0.72)
than for dengue (0.88). Its 80\% intervals cover 81\% of outcomes overall, against 78\% for the baseline, so
the bands mean what they say.

\begin{teammate}{Samira (AN-5): confirm}
Horizons 2 and 3 are in \texttt{backtest\_metrics.csv}. Add the WIS decomposition (for dengue, most of the model's WIS
is underprediction, which matters for stock planning). Replace with the v1.0 numbers.
\end{teammate}

\subsection{Allocation against the population baseline}
\begin{teammate}{Rafi}
In the SF=0.25 run the national net need is zero for beds, IV saline, NS1 kits and ORS, and 8.8 courses of
antimalarials, because the generated stock exceeds the forecast demand. The LP's ``50\% of need met against 1\%
for the population baseline'' is therefore about antimalarials only. Fix the demand or stock assumptions, then
report share of need met and the worst-covered district at five budgets, and the efficiency and equity curve.
\end{teammate}

\subsection{Scalability and Spark optimisations}
\texttt{scripts/scaling.sh} builds each scale factor from a fresh clone and times every stage
(Table~\ref{tab:scaling}). The laptop has an Intel Core i5-8350U (4 cores, 8 threads) and 15.4\,GB of RAM;
Spark runs as \texttt{local[4]} with a 5\,GB driver. SF=1 finishes in 28.3 minutes, inside the 45-minute
gate. Silver grows with the row count and is the slowest stage at SF=1; analytics stays near 7 to 8 minutes at
every scale, because it works on weekly series per district, not on rows. A first SF=1 attempt failed in silver
when Spark's spill filled the small RAM-backed \texttt{/tmp}; spill now goes to a disk folder.

\begin{table}[htbp]
\caption{Minutes per stage from a fresh clone}
\label{tab:scaling}
\centering\small
\setlength{\tabcolsep}{4pt}
\begin{tabular}{@{}lrrr@{}}
\toprule
 & SF=0.25 & SF=0.5 & SF=1 \\
\midrule
Encounter rows (M) & 8.5 & 17.1 & 34.1 \\
Generate & 0.2 & 0.4 & 0.6 \\
Bronze & 1.9 & 2.8 & 4.3 \\
Silver & 3.8 & 6.7 & 11.1 \\
Gold & 1.8 & 2.3 & 3.5 \\
Analytics & 6.9 & 7.1 & 8.0 \\
Serve & 0.7 & 1.1 & 0.9 \\
\midrule
Total & 15.4 & 20.4 & 28.3 \\
Raw / silver (GB) & 1.2 / 0.8 & 2.2 / 1.5 & 4.1 / 2.9 \\
\bottomrule
\end{tabular}
\end{table}

One optimisation is measured. The silver job evaluated the 14 encounter quality rules, including the
duplicate-ID rule that sorts the whole load, once for every output (quarantine, its count, the quality report
and the clean MERGE). It now evaluates them once and persists the result. At SF=0.25 silver's Spark time fell
from 325\,s to 214\,s, and every silver table kept the same row count and row hash
(\texttt{silver\_optimisation.json}).

\begin{teammate}{Samira (ENG-3)}
Spark UI screenshots from the SF=1 event log (kept outside the repository; it holds the padded HMAC key), the
bottleneck stage, and a second optimisation with before and after numbers (candidates: AQE skew handling or
salting on the Dhaka load; backtest parallelism, 365\,s at SF=1 on 4 of 8 threads).
\end{teammate}
